\section{Introduction}
\label{sec:introduction}

Le développement de systèmes logiciels complexes nécessite aujourd'hui des environnements capables d'accompagner les différentes étapes de leur conception, de leur validation et de leur évolution. Dans le domaine des systèmes à composants, ces environnements doivent non seulement permettre de modéliser des architectures de plus en plus élaborées, mais également de proposer des outils capables d'évoluer avec les besoins des utilisateurs et des projets de recherche. \\

Ce rapport s'inscrit dans le cadre du stage de Master 2 Ingénierie Système et Logiciel. Réalisé au sein de l'équipe VESONTIO de l'institut FEMTO-ST, il contribue au développement du langage CHIPS et de son environnement logiciel. Les travaux effectués durant ce projet portent sur l'évolution de plusieurs composants de la chaîne de compilation ainsi que sur l'enrichissement des fonctionnalités offertes par le langage. \\

Les objectifs poursuivis durant ce projet sont multiples. Ils consistent d'une part à améliorer certains éléments de l'infrastructure logicielle afin de faciliter l'évolution du langage, et d'autre part à intégrer de nouveaux composants réutilisables destinés à étendre les possibilités offertes aux utilisateurs. Enfin, une étude expérimentale est menée afin d'évaluer les performances des différentes solutions développées et leur impact sur la chaîne de compilation. \\

Ce rapport est organisé en six sections. \\

La section \ref{sec:pres_entreprise} présente la structure d'acceuil dans laquelle s'inscrit ce stage. \\

La section \ref{sec:etats_de_art} expose le contexte scientifique nécessaire à la compréhension des travaux 
réalisés. \\

La section \ref{sec:parseur_XMI} décrit les évolutions apportées au compilateur du langage CHIPS. \\ 

La section \ref{sec:outils_theorie_controle} présente les nouveaux composants développés ainsi que leur 
évaluation expérimentale. \\

Enfin, la section \ref{sec:conclusion} conclut ce rapport en présentant un bilan des travaux réalisés 
et plusieurs perspectives d'évolution.


\FloatBarrier

\newpage